\chapter{Objetivos}

\section{Objetivo general}

Desarrollar un modelo computacional de un acelerador
lineal de protones con estructura tipo Alvarez, mediante
la implementación de un modelo de aceleración por
radiofrecuencia, con el propósito de estudiar el diseño
de las celdas y analizar la dinámica longitudinal y la
estabilidad de fase de un paquete de partículas.

\section{Objetivos específicos}

\begin{enumerate}

    \item Implementar el modelo de aceleración por cavidades
    de radiofrecuencia, considerando la ganancia de energía
    cinética por hueco y la energía acumulada después de
    múltiples etapas de aceleración:

    \begin{equation}
        \Delta T_i = qV_i\cos\phi_i,
    \end{equation}

    \begin{equation}
        T_N = T_0 + \sum_{i=1}^{N} qV_i\cos\phi_i.
    \end{equation}

    \item Determinar las longitudes de las celdas del
    acelerador a partir de la condición de sincronismo
    correspondiente a la partícula síncrona.

    \item Analizar la influencia de la fase síncrona y
    de los errores de sincronización de la radiofrecuencia
    sobre la dinámica longitudinal y el comportamiento
    de un paquete de partículas.

    \item Contrastar los resultados de las simulaciones
    de las oscilaciones de sincrotrón con las predicciones
    de la teoría lineal y comparar el comportamiento
    longitudinal de protones y electrones.

\end{enumerate}

